\documentclass[10pt,letterpaper]{article}
\usepackage{resumen}
\AtBeginDocument{\shorthandoff{<>}}

% Ajustes locales de espaciado (más aire entre bloques)
\setlength{\parskip}{4pt plus 1pt}
\tcbset{clave/.append style={top=5pt,bottom=5pt,before skip=7pt,after skip=7pt},
        apoyo/.append style={top=4pt,bottom=4pt,before skip=8pt,after skip=8pt},
        nota/.append style={top=4pt,bottom=4pt,before skip=8pt,after skip=8pt},
        alerta/.append style={top=4pt,bottom=4pt,before skip=8pt,after skip=8pt}}
\titlespacing*{\section}{0pt}{10pt plus 2pt}{5pt}

\begin{document}
\titulo{Cinética química II}{Ecuaciones integradas, vida media y dependencia con la temperatura}

%------------------------------------------------------------------
\section{Ecuaciones integradas y tiempo de vida media}
Integrar la ley de velocidad permite conocer la concentración en función del tiempo. Para A $\to$ productos con orden $n$, se separan las variables y se integra entre $[\mathrm{A}]_0$ (en $t=0$) y $[\mathrm{A}]$ (en $t$):
\begin{clave}
$\displaystyle r=-\frac{d[\mathrm{A}]}{dt}=k\,[\mathrm{A}]^{n}\quad\Longrightarrow\quad \int_{[\mathrm{A}]_0}^{[\mathrm{A}]}\frac{d[\mathrm{A}]}{[\mathrm{A}]^{n}}=-k\int_0^{t}dt$
\end{clave}

La vida media $t_{1/2}$ es el tiempo que tarda $[\mathrm{A}]$ en reducirse a $[\mathrm{A}]_0/2$; al imponer esa condición en cada ecuación integrada resulta la última columna.

\begin{center}\small
\renewcommand{\arraystretch}{2.0}
\begin{tabular}{@{}c c c c c c@{}}
\toprule
Orden & Ecuación integrada & Recta contra $t$ & Pendiente & Unidades de $k$ & $t_{1/2}$\\
\midrule
0 & $[\mathrm{A}]=[\mathrm{A}]_0-kt$ & $[\mathrm{A}]$ & $-k$ & \si{\molar\per\second} & $\dfrac{[\mathrm{A}]_0}{2k}$\\
1 & $\ln[\mathrm{A}]=\ln[\mathrm{A}]_0-kt$ & $\ln[\mathrm{A}]$ & $-k$ & \si{\per\second} & $\dfrac{\ln 2}{k}$\\
2 & $\dfrac{1}{[\mathrm{A}]}-\dfrac{1}{[\mathrm{A}]_0}=kt$ & $1/[\mathrm{A}]$ & $+k$ & \si{\per\molar\per\second} & $\dfrac{1}{k[\mathrm{A}]_0}$\\
\bottomrule
\end{tabular}
\end{center}

En forma explícita, $[\mathrm{A}]=[\mathrm{A}]_0\,e^{-kt}$ en primer orden (con $\ln 2=0{,}693$) y $[\mathrm{A}]=[\mathrm{A}]_0/(1+kt[\mathrm{A}]_0)$ en segundo. La vida media crece con $[\mathrm{A}]_0$ en orden 0, no depende de ella en orden 1 y disminuye con ella en orden 2.

\vspace{4pt}
\begin{minipage}[t]{0.50\linewidth}\vspace{0pt}
\begin{tikzpicture}
\begin{axis}[resumen,width=\linewidth,height=3.9cm,xlabel={$t$ (en unidades de la primera $t_{1/2}$)},ylabel={$[\mathrm{A}]/[\mathrm{A}]_0$},
  xmin=0,xmax=4,ymin=0,ymax=1.05,xtick={0,1,2,3,4},ytick={0,0.5,1},legend pos=north east,legend style={font=\scriptsize,draw=none,fill=white,fill opacity=0.8,text opacity=1}]
\addplot[acc,domain=0:2,samples=2]{1-x/2};
\addplot[teal,domain=0:4,samples=60]{2^(-x)};
\addplot[bad,domain=0:4,samples=60]{1/(1+x)};
\addplot[dashed,gray,line width=0.6pt] coordinates {(0,0.5)(4,0.5)};
\legend{orden 0,orden 1,orden 2}
\end{axis}
\end{tikzpicture}
\end{minipage}\hfill
\begin{minipage}[t]{0.46\linewidth}\vspace{4pt}\small\raggedright
\textbf{Vidas medias sucesivas.} Si cada intervalo parte de la mitad de la concentración anterior, en orden 0 cada vida media dura la mitad de la anterior; en orden 1 todas duran lo mismo y queda $1/2$, $1/4$, $1/8$ del reactivo; en orden 2 cada una dura el doble.
\vspace{8pt}

\textbf{Ejemplo, orden 1.} Una dosis de $^{99\mathrm{m}}$Tc ($t_{1/2}=\SI{6,01}{\hour}$) tiene $A_0=\SI{740}{\mega\becquerel}$. Con $k=\ln 2/t_{1/2}=\SI{0,1153}{\per\hour}$, a las \SI{6,5}{\hour} queda $A_0e^{-kt}=\SI{350}{\mega\becquerel}$.
\end{minipage}

%------------------------------------------------------------------
\section{Determinación del orden}
El orden se obtiene por velocidades iniciales (Cinética I), por ajuste de las ecuaciones integradas o por el tiempo de vida media. Las ecuaciones integradas requieren $[\mathrm{A}]$ medida en el tiempo en un experimento; la vida media, $t_{1/2}$ para distintas $[\mathrm{A}]_0$.

\subsection*{Ecuaciones integradas}
Se linealizan los datos como $[\mathrm{A}]$, $\ln[\mathrm{A}]$ y $1/[\mathrm{A}]$ contra $t$, se ajusta una recta a cada serie y se adopta el orden de la de mayor $R^{2}$, siempre que los residuos no muestren curvatura sistemática. Un fármaco en disolución, con 8 mediciones entre 0 y \SI{16}{\hour}, da $R^{2}=0{,}9832$ (orden 0), $0{,}9966$ (orden 1) y $0{,}9891$ (orden 2); se adopta el orden 1, y la pendiente de $\ln[\mathrm{A}]$, $-0{,}0441$, entrega $k=\SI{0,0441}{\per\hour}$.

\clearpage
\subsection*{Método del tiempo de vida media}
Con $[\mathrm{A}]=[\mathrm{A}]_0/2$ y $t=t_{1/2}$ en la ecuación integrada general $\dfrac{1}{n-1}\Big(\dfrac{1}{[\mathrm{A}]^{n-1}}-\dfrac{1}{[\mathrm{A}]_0^{\,n-1}}\Big)=kt$ se obtiene
\begin{clave}
$\displaystyle t_{1/2}=\frac{2^{\,n-1}-1}{(n-1)\,k\,[\mathrm{A}]_0^{\,n-1}}\qquad\Longrightarrow\qquad
\log t_{1/2}=\log\frac{2^{\,n-1}-1}{(n-1)\,k}-(n-1)\,\log[\mathrm{A}]_0$
\end{clave}
El gráfico de $\log t_{1/2}$ contra $\log[\mathrm{A}]_0$ es una recta de pendiente $-(n-1)$, de donde $n=1-\text{pendiente}$; el intercepto (coeficiente de posición) permite despejar $k$ una vez conocido $n$. El método no es aplicable a $n=1$, pues entonces $t_{1/2}=\ln 2/k$ no depende de $[\mathrm{A}]_0$ y la pendiente es nula.

\vspace{4pt}
\begin{minipage}[t]{0.50\linewidth}\vspace{0pt}
\begin{tikzpicture}
\begin{axis}[resumen,width=\linewidth,height=3.6cm,xlabel={$\log[\mathrm{A}]_0$},ylabel={$\log t_{1/2}$},
  xmin=-1.15,xmax=0.05,ymin=1.4,ymax=2.5,xtick={-1,-0.8,-0.6,-0.4,-0.2,0},ytick={1.5,2,2.5}]
\addplot[only marks,mark=*,acc,mark size=2pt] coordinates {(-1.0,2.4014)(-0.69897,2.0934)(-0.39794,1.7993)(-0.09691,1.4983)};
\addplot[teal,domain=-1.15:0.05,samples=2]{1.4009-0.9977*x};
\node[font=\scriptsize,anchor=south west] at (axis cs:-0.75,2.30) {pendiente $=-0{,}998$};
\end{axis}
\end{tikzpicture}
\end{minipage}\hfill
\begin{minipage}[t]{0.46\linewidth}\vspace{2pt}\small\raggedright
\textbf{Ejemplo: peróxido en disolución.} Con $[\mathrm{A}]_0=0{,}10;\ 0{,}20;\ 0{,}40;\ \SI{0,80}{\molar}$ se midió $t_{1/2}=252;\ 124;\ 63{,}0;\ \SI{31,5}{\second}$. La recta de ajuste tiene pendiente $-0{,}998$ e intercepto $1{,}401$:
\begin{align*}
n&=1-(-0{,}998)=1{,}998\approx2\\
\log\tfrac{1}{k}&=1{,}401\ \Rightarrow\ k=10^{-1{,}401}\\
&=\SI{0,040}{\per\molar\per\second}
\end{align*}
\vspace{-8pt}
\end{minipage}

\vspace{2pt}
{\small\textit{Verificación de sanidad.} En orden 2, al duplicar $[\mathrm{A}]_0$ la vida media se reduce a la mitad ($252\to\SI{124}{\second}$), coherente con $n\approx2$.}

%------------------------------------------------------------------
\section{Dependencia de $k$ con la temperatura}
Según la teoría de colisiones, reaccionan las moléculas que chocan con orientación adecuada y energía suficiente, y la energía de activación $E_a$ es la energía umbral. La relación de Arrhenius, en sus formas exponencial, lineal y de dos temperaturas, es
\begin{clave}
$\displaystyle k=A\,e^{-E_a/RT}\qquad\qquad \ln k=\ln A-\frac{E_a}{R}\,\frac{1}{T}\qquad\qquad \ln\frac{k_2}{k_1}=-\frac{E_a}{R}\Big(\frac{1}{T_2}-\frac{1}{T_1}\Big)$
\end{clave}
Aquí $A$ es el factor preexponencial o de colisión, con las unidades de $k$, y $E_a$ la energía de activación, en \si{\joule\per\mole} si $R=\SI{8,314}{\joule\per\mole\per\kelvin}$ y $T$ está en kelvin. En el gráfico de $\ln k$ contra $1/T$ la pendiente es $-E_a/R$ y el intercepto, $\ln A$. Un catalizador ofrece una ruta con menor $E_a$ y no modifica la diferencia de energía entre reactivos y productos.

\vspace{2pt}
\begin{minipage}[t]{0.49\linewidth}\vspace{0pt}
\begin{tikzpicture}
\begin{axis}[resumen,width=\linewidth,height=3.7cm,xlabel={$1000/T$ (\si{\per\kelvin})},ylabel={$\ln k$},
  xmin=2.95,xmax=3.42,ymin=-8.4,ymax=-3.4,xtick={3,3.1,3.2,3.3,3.4},ytick={-8,-7,-6,-5,-4}]
\addplot[only marks,mark=*,acc,mark size=2pt] coordinates {(3.1933,-5.8091)(3.0945,-4.8728)(3.0017,-3.9737)};
\addplot[teal,dashed,domain=2.95:3.42,samples=2]{24.7548-9.57211*x};
\addplot[only marks,mark=*,bad,mark size=2.4pt] coordinates {(3.3540,-7.3500)};
\node[font=\scriptsize,bad,anchor=north] at (axis cs:3.354,-7.5) {25\,\textdegree C};
\end{axis}
\end{tikzpicture}
\end{minipage}\hfill
\begin{minipage}[t]{0.47\linewidth}\vspace{0pt}
\begin{tikzpicture}[x=1cm,y=1cm,font=\scriptsize,scale=0.8,transform shape]
\draw[->,line width=0.6pt] (0,0) -- (0,3.6) node[above] {energía};
\draw[->,line width=0.6pt] (0,0) -- (6.4,0) node[right] {avance};
\draw[acc,line width=1.3pt] (0.3,1.0) .. controls (1.4,1.0) and (1.8,3.3) .. (3.1,3.3) .. controls (4.4,3.3) and (4.8,0.4) .. (6.0,0.4);
\draw[teal,dashed,line width=1.2pt] (0.3,1.0) .. controls (1.4,1.0) and (1.8,2.1) .. (3.1,2.1) .. controls (4.4,2.1) and (4.8,0.4) .. (6.0,0.4);
\draw[gray,dotted] (0.3,1.0) -- (3.7,1.0);
\draw[<->,bad,line width=0.7pt] (3.55,1.0) -- (3.55,3.3) node[midway,right] {$E_a$};
\draw[<->,teal,line width=0.7pt] (2.6,1.0) -- (2.6,2.1) node[midway,right] {$E_{a,\mathrm{cat}}$};
\node[anchor=north west] at (0.3,1.0) {Reactivos};
\node[anchor=north east] at (6.2,0.3) {Productos};
\draw[acc,line width=1.3pt] (4.3,3.55) -- (4.8,3.55) node[right,black] {sin catalizador};
\draw[teal,dashed,line width=1.2pt] (4.3,3.2) -- (4.8,3.2) node[right,black] {con catalizador};
\end{tikzpicture}
\end{minipage}

\vspace{2pt}
\begin{minipage}[t]{0.49\linewidth}\vspace{0pt}\small\raggedright
\textbf{Ejemplo: estabilidad acelerada.} La degradación de una solución oral es de primer orden, con $k=0{,}00300;\ 0{,}00765;\ \SI{0,0188}{\per\day}$ a $40;\ 50;\ \SI{60}{\celsius}$. Se aplica la forma lineal y se extrapola a \SI{25}{\celsius}.
\end{minipage}\hfill
\begin{minipage}[t]{0.47\linewidth}\vspace{-4pt}\small
\begin{align*}
\ln k&=24{,}755-\frac{\SI{9572}{\kelvin}}{T}\\
E_a&=(\SI{9572}{\kelvin})\,R=\SI{79,6}{\kilo\joule\per\mole}\\
k_{25}&=\SI{6,42e-4}{\per\day}\ \Rightarrow\ t_{90}=\frac{\ln(1/0{,}9)}{k}=\SI{164}{\day}
\end{align*}
\vspace{-8pt}
\end{minipage}

\begin{alerta}\small
\textbf{Errores frecuentes.} Usar $T$ en \si{\celsius}; combinar $E_a$ en \si{\kilo\joule\per\mole} con $R$ en \si{\joule\per\mole\per\kelvin}; extrapolar fuera del intervalo de validez.
\end{alerta}
\end{document}
